\begin{claim}{D.78}{Proved}{The other negative fixed subgroup and its actual pro-odd kernel}
Retain the two-section-class situation of D.75:
\[
 F=\widehat F(x,y),\quad c=[x,y],\quad
 E=F\rtimes\langle v\rangle,\quad vcv^{-1}=c^{-1},\quad
 gvg^{-1}=cv.
\]
The transporter $g$ is the unique original D.75 transporter. Put
$S=\langle v\rangle$, let $T=\langle w\rangle$ be an actual representative
of the other section class, and put
\[
 B=C_F(w),\qquad N_E(T)=B\times T,\qquad k=\mathbb F_2.
\]
D.59 gives $\dim_kH^1(B,k)=1$. Retain the actual D.60 quotient
$\chi:E\to\widehat{\mathbb Z}$, normalized by $\chi(g)=1$.
Its restriction to $B$ is onto. Choose $a\in B$ with
$[a]_{\mathrm{ab}}=[g]_{\mathrm{ab}}$, as in D.60, so $\chi(a)=1$.
Define the actual closed groups
\[
 A=\overline{\langle a\rangle},\qquad
 K=\ker(\chi|_B),\qquad
 O=\ker(B\xrightarrow{\chi}\widehat{\mathbb Z}\to\mathbb Z_2).
\]

Every open subgroup $B_0$ of the full $B$ satisfies
\[
                         \dim_k H^1(B_0,k)=1.                 \tag{D.78.1}
\]
The actual subgroup $A_2=\overline{\langle a^{e_2}\rangle}$ is a
Sylow pro-$2$ subgroup of $B$, and the restriction of $\chi_2$ to it
is an isomorphism onto $\mathbb Z_2$. Both $O$ and $K$ are pro-odd.
In particular the full original fixed group has the exact inverse coordinates
\[
 \begin{aligned}
 K\rtimes\widehat{\mathbb Z}&\longrightarrow B,
       &(k,t)&\longmapsto k a^t,\\
 B&\longrightarrow K\rtimes\widehat{\mathbb Z},
       &b&\longmapsto(ba^{-\chi(b)},\chi(b)).
 \end{aligned}                                                \tag{D.78.2}
\]
The action is conjugation by the specified $a^t$.
For every finite continuous $2$-primary $B$-module $Q$, the actual inflation map is
an isomorphism
\[
       H^n(\widehat{\mathbb Z},Q^K)
                       \xrightarrow{\ \sim\ }H^n(B,Q).       \tag{D.78.3}
\]
Under $A\xrightarrow{\chi}\widehat{\mathbb Z}$, restriction to $A$
is precisely the coefficient-inclusion map
\[
 H^n(\widehat{\mathbb Z},Q^K)
                \longrightarrow H^n(\widehat{\mathbb Z},Q).   \tag{D.78.4}
\]
The action on $Q$ in the right-hand group is explicitly
$t\cdot q=a^t q$; on $Q^K$ it is the quotient action of $B/K$.
These actions agree on the indicated submodule.
For every finite continuous $E$-set $X$, put $Q_X=k[X^T]$ with its
actual $B$-action. The proof supplies perfect finite-coefficient pairings
and the exact comparison $\mathcal U_T(k[X])\cong
H^0(B,Q_X)\oplus H^1(B,Q_X)$. The full recognition $B=A$ is equivalent
to surjectivity for every such $X$ of the specified inclusion
$Q_X^B\to Q_X^A$. Every nonidentity element of $K$ has the explicit
finite-coset coefficient witness given below.
Triviality of $K$ and that invariant-surjectivity assertion remain
unproved. The original cusp $c$, transporter $g$, and D.75.5 remain
unchanged.

\end{claim}
\begin{proof}
\emph{The actual perfect localized duality.}
Use the original cusp normalizer $D=N_E(\overline{\langle c\rangle})$
and the exact augmentation kernel and coefficient modules of D.75.7--13.
For a finite continuous $E$-module $M$, write
\[
 \begin{aligned}
 \mathcal A_M&=H^\bullet(E,M)[u^{-1}]_0,\\
 \mathcal B_M&=H^\bullet(D,M)[u^{-1}]_0,\\
 \mathcal R_M&=H^\bullet(E,D;M)[u^{-1}]_0,
 \end{aligned}
\]
where $u:E\to C_2=k$ is the original quotient character. The maps
\[
 \mathcal A_M\xrightarrow{r}\mathcal B_M
       \xrightarrow{u^{-1}\delta}\mathcal R_M
       \xrightarrow{j}\mathcal A_M                         \tag{D.78.5}
\]
are the actual localized restriction, connecting and relative-to-absolute
maps. They are linear over the localized coefficient ring $\mathcal A_k$.

The D.75.14 transgression calculation gives more than equality of ranks.
Put $U_M=H^0(F,M)$ and $V_M=H^1(F,M)$. In positive quotient columns the
absolute rows are $\overline U_M,\overline V_M$, whereas the shifted
relative rows for $M^\vee$ are
$\overline{V_M^\vee},\overline{U_M^\vee}$, with the kernel/image cyclic
pairings of D.74.19. The only possible differentials are transpose maps
\[
 d:\overline V_M\to\overline U_M,
 \qquad d':\overline{U_M^\vee}\to\overline{V_M^\vee}.
\]
Consequently the pairings on the two surviving pairs of rows are perfect:
\[
 \operatorname{coker}d\ \times\ \ker d'\longrightarrow k,
 \qquad
 \ker d\ \times\ \operatorname{coker}d'\longrightarrow k.   \tag{D.78.6}
\]
These are the actual pairings induced by evaluation and the fibre
surface-pair fundamental class, not chosen vector-space identifications.
For sufficiently large absolute degree $n$ and relative degree $m$, the
surviving filtration sequences are exactly
\[
 \begin{gathered}
 0\longrightarrow\operatorname{coker}d\longrightarrow H^n(E,M)
      \longrightarrow\ker d\longrightarrow0,\\
 0\longrightarrow\operatorname{coker}d'\longrightarrow H^m(E,D;M^\vee)
      \longrightarrow\ker d'\longrightarrow0.
 \end{gathered}
\]
The first subspace has quotient-column filtration $n$, and the second
has filtration $m-1$ in the shifted relative complex.

Here is the required trace on the target. For trivial coefficient $k$,
the absolute transgression is zero: the actual section $v$ makes
inflation from $C_2$ injective in every degree. Its transpose in the
shifted relative sequence is therefore zero as well. The top fibre row
$H^1(F,\mathscr H_k)=H^2(F,\langle c\rangle;k)=k$ survives.
For sufficiently large $n$ its quotient edge map supplies
\[
 H^n(E,D;k)\longrightarrow
 E_\infty^{n-2,1}=H^{n-2}(C_2,k)=k u^{n-2}\longrightarrow k,
                                                               \tag{D.78.7}
\]
where the last map reads the coefficient of the specified power of $u$.
These maps commute with a further multiplication by $u$. Thus they give
a specified trace $\operatorname{tr}:\mathcal R_k\to k$.

Cup product followed by this trace gives
\[
 \mathcal A_M\times\mathcal R_{M^\vee}\longrightarrow k.       \tag{D.78.8}
\]
Its associated two-row pairings are exactly (D.78.6), in opposite filtration
positions. The product of the two filtration subspaces just displayed has
quotient filtration at least $n+m-1$, whereas the trace reads column
$n+m-2$; that product therefore has trace zero. An absolute vector with
nonzero image in $\ker d$ pairs nontrivially with the relative subspace
$\operatorname{coker}d'$. A nonzero absolute vector in
$\operatorname{coker}d$ pairs nontrivially with a relative vector mapping
to $\ker d'$. These statements follow from the two perfect pairings in
(D.78.6), and the same argument applies with the sides reversed. This proves
nondegeneracy on the full finite filtered spaces, not only on their
associated graded spaces. Passing through the specified $u$ maps proves
that (D.78.8) is perfect.
The construction is natural in the coefficient maps and balanced over
$\mathcal A_k$, by the original cup-product and connecting-map identities.

\emph{The other Brown factor and its exact relative lift.}
D.75 proves $C_F(v)=1$. Brown restriction therefore gives the actual
coefficient decomposition
\[
 \mathcal A_M\xrightarrow{\ \sim\ }
 \mathcal V_S(M)\oplus\mathcal U_T(M),\qquad
 \begin{cases}
 \mathcal V_S(M)=H^\bullet(S,M)[u^{-1}]_0,\\
 \mathcal U_T(M)=H^\bullet(B\times T,M)[u^{-1}]_0.
 \end{cases}                                                  \tag{D.78.9}
\]
The two cusp factors are transported by $1$ and the original $g$, with
exactly the coefficient transport of D.75.18. The restriction in (D.78.5)
is consequently
\[
 \mathcal V_S(M)\oplus\mathcal U_T(M)
            \longrightarrow\mathcal V_S(M)\oplus\mathcal V_S(M),
              \qquad(z,t)\longmapsto(z,z).                    \tag{D.78.10}
\]
In particular the other factor has zero boundary restriction for every
finite module, including every finite permutation module.

For coefficient $k$, let $e\in\mathcal A_k$ be the Brown inverse of
$(0,1)$, the identity in the other normalizer factor and zero in the cusp
factor. It is a specified idempotent. Its restriction to both cusp factors
is zero, so $e\mathcal B_M=0$. Applying this idempotent to (D.78.5) gives
an actual isomorphism
\[
            j:e\mathcal R_M\xrightarrow{\ \sim\ }
                         e\mathcal A_M=\mathcal U_T(M).       \tag{D.78.11}
\]
Its inverse is unambiguous: lift an element through $j$, then multiply
its lift by $e$. Two lifts differ by a connecting image, which $e$ kills.

The pairing (D.78.8) splits orthogonally under $e$ and $1-e$; its restriction
to $e\mathcal A_M\times e\mathcal R_{M^\vee}$ is perfect. Transporting it
through (D.78.11) gives a perfect pairing
\[
 \mathcal U_T(M)\times\mathcal U_T(M^\vee)\longrightarrow k,
 \qquad(z,z')\longmapsto\tau\bigl(\operatorname{ev}_*(z\smile z')\bigr),
                                                               \tag{D.78.12}
\]
where $\tau=\operatorname{tr}\,j^{-1}:\mathcal U_T(k)\to k$ uses
(D.78.11). The formula follows because $j$ commutes with cup multiplication
by an absolute class; it therefore retains the actual normalizer cup
product and the actual coefficient evaluation.

\emph{Arbitrary finite permutation coefficients.}
Let $X$ be any finite continuous $E$-set and $M=k[X]$, with its invariant
basis pairing. The $T$-fixed points give a specified $B\times T$-submodule
\[
                         Q_X=k[X^T]\subseteq M.              \tag{D.78.13}
\]
Projection onto these basis elements is an inverse on this summand.
The complementary summand consists of the two-point $T$-orbits and is
free as a $k[T]$-module. It is $B$-stable because $B$ centralizes $T$.
Its positive $T$-cohomology vanishes. The extension spectral sequence
for $T\to B\times T\to B$, with $\operatorname{cd}_2B\leq1$, shows
that this complementary summand contributes no cohomology in degree
at least two. Thus inclusion of $Q_X$ induces the exact localized
coefficient comparison.

The action of $T$ on $Q_X$ is trivial. The direct-product resolution,
or the two-row complex with zero $T$ differential, gives the actual
isomorphism
\[
 \begin{aligned}
 H^0(B,Q_X)\oplus H^1(B,Q_X)&\xrightarrow{\ \sim\ }\mathcal U_T(M),\\
 (q,z)&\longmapsto q+u^{-1}\operatorname{inf}_B(z).
 \end{aligned}                                                \tag{D.78.14}
\]
At a positive degree $n\geq2$, this formula means
$q u^n+\operatorname{inf}_B(z)u^{n-1}$, followed by the coefficient
inclusion (D.78.13). Projection to the fixed-point summand and the direct
product cochain decomposition supply its inverse.
Restriction to the literal subgroup $T$ is exactly
\[
 H^0(B,Q_X)\oplus H^1(B,Q_X)\longrightarrow Q_X,
                         \qquad(q,z)\longmapsto q.           \tag{D.78.15}
\]
Thus its image is $Q_X^B$, with precisely that subgroup action retained.

The $H^1(B,Q_X)$ summand in (D.78.14) is totally isotropic for (D.78.12):
the product of two such summands is a coefficient image of
$H^2(B,Q_X\otimes Q_X)=0$. Since (D.78.12) is perfect on the finite
self-dual module $M$, it follows that
\[
                 \dim H^1(B,Q_X)\leq\dim H^0(B,Q_X).         \tag{D.78.16}
\]
Decompose the actual finite $B$-set $X^T$ into its orbits, with
stabilizers $B_1,\ldots,B_r$. The usual finite-coset Shapiro maps give
\[
 H^0(B,Q_X)=k^r,\qquad
 H^1(B,Q_X)=\bigoplus_{i=1}^rH^1(B_i,k).                       \tag{D.78.17}
\]
Every $B_i$ is open in $B$, and $\chi(B_i)$ has finite index in
$\widehat{\mathbb Z}$ because $\chi(B)=\widehat{\mathbb Z}$.
Write this image as $d_i\widehat{\mathbb Z}$, with $d_i$ a positive
integer. The actual map
\[
       B_i\xrightarrow{\chi}d_i\widehat{\mathbb Z}
        \xrightarrow{\div d_i}\widehat{\mathbb Z}
          \longrightarrow\mathbb F_2                        \tag{D.78.18}
\]
is a nonzero continuous character. Division here is the continuous
inverse of multiplication by $d_i$ onto its indicated image.
Thus each summand in (D.78.17) has dimension at least one. Together with
(D.78.16) this proves equality and dimension exactly one for each stabilizer.

Every open subgroup $B_0\leq B$ occurs as such a stabilizer. Choose an
open normal $N\triangleleft E$ contained in $F$ with
$N\cap B\subseteq B_0$. The group $U=B_0TN$ is open in $E$, since
$B_0T$ is a subgroup and $N$ is normal. The quotient label $u$ and
$N\subset F$ show $U\cap B=B_0$. In the actual finite $E$-set
$X=E/U$, the identity coset is $T$-fixed and its $B$-stabilizer is $B_0$.
This proves (D.78.1), with no passage to a maximal pro-$2$ quotient.

The two summands in (D.78.14) now have equal dimension. The mixed part
of the perfect pairing (D.78.12) is therefore perfect as well. Indeed
a vector of $H^1$ pairing trivially with $H^0$ pairs trivially with
both summands and must vanish, and the dimensions agree. Hence the
actual cup pairings, with the degree-correct trace
$\tau_1(z)=\tau(u^{-1}\operatorname{inf}_Bz)$, are
\[
 H^0(B,Q_X)\times H^1(B,Q_X^\vee)
     \longrightarrow H^1(B,k)\xrightarrow{\tau_1}k            \tag{D.78.19}
\]
and their reversed versions are perfect. The fixed-point basis pairing
identifies $Q_X^\vee$ with $Q_X$. On the one-dimensional $H^1(B,k)$,
$\tau_1$ is nonzero: otherwise that square-zero summand would lie in the
radical of (D.78.12) for $M=k$. Thus $\tau_1$ takes the original nonzero class
$\chi|_B\bmod2$ to one. This records the original finite-coefficient duality, including the
actual stabilizers and coefficient actions.

\emph{The full pro-odd kernel consequence.}
The map $t\mapsto a^t$ is an actual continuous section of $\chi|_B$.
In particular $A_2=\overline{\langle a^{e_2}\rangle}$ is a copy of
$\mathbb Z_2$ mapping isomorphically onto the dyadic quotient. Choose a
Sylow pro-$2$ subgroup $P$ of $B$ containing $A_2$.
There cannot be two linearly independent continuous characters
$P\to k$. Their joint map factors through $P/(P\cap N_0)$ for some
open normal $N_0\triangleleft B$. It therefore extends to the open
subgroup $PN_0$ by the quotient isomorphism
$PN_0/N_0=P/(P\cap N_0)$, contradicting (D.78.1).

It follows that $P$ is procyclic. One elementary verification is that
every nontrivial finite $2$-group quotient has at most one independent
character to $C_2$, hence a unique maximal subgroup; any element outside
that maximal subgroup generates the finite group. Thus every finite
quotient is cyclic, and compatible generators give a procyclic inverse
limit. Since $P$ maps onto $\mathbb Z_2$, it is $\mathbb Z_2$, and that
surjection is an isomorphism. Its subgroup $A_2$ already maps onto the
same quotient, so $P=A_2$.

Every Sylow pro-$2$ subgroup of $B$ is conjugate to $A_2$.
The homomorphism $\chi_2$ is invariant under conjugation and is injective
on all these subgroups. Hence $O=\ker\chi_2$ has trivial Sylow pro-$2$
subgroup and all its finite quotients have odd order. It is pro-odd.
The actual subgroup $K=\ker(\chi|_B)$ lies in $O$, so is pro-odd too.
The multiplication and inverse formulas (D.78.2) follow directly from
the chosen section $a^t$ and the original character $\chi$.

\emph{The surviving coefficient map and the original scope.}
For any finite continuous $2$-primary $B$-module $Q$, average over the
finite odd image of $K$ on $Q$. Its odd order is invertible on $Q$. The averaging map
\[
       e_K:Q\longrightarrow Q^K,
       \qquad q\longmapsto\frac1{|\overline K|}
                               \sum_{h\in\overline K}hq       \tag{D.78.20}
\]
is independent of the chosen finite image through which the action
factors. It commutes with $B$, since $K$ is normal. It makes the
invariants functor exact; higher $2$-primary cohomology of $K$ vanishes.
The extension spectral sequence for the actual sequence
$K\to B\xrightarrow{\chi}\widehat{\mathbb Z}$ therefore gives (D.78.3).
Restriction to the specified section $A$ commutes with inflation and
is induced by the inclusion $Q^K\hookrightarrow Q$, proving (D.78.4).
Evaluation, coefficient maps and cup products commute with this comparison.

For $Q=Q_X=k[X^T]$, the invariant space is concretely
$k[K\backslash X^T]$, mapped to $Q_X$ by its orbit sums. Every such orbit
has odd size, so the original basis pairing restricts to the ordinary
basis pairing on these sums. The action of the quotient is the action
of the specified powers of $a$ on the $K$-orbits. This is the actual
finite-coefficient comparison; it retains the action of $K$ through
its orbit partition.

There is an exact finite-permutation detection of the surviving kernel.
The equality $B=A$ is equivalent to surjectivity, for every finite
continuous $E$-set $X$, of the actual restriction map
\[
       Q_X^B\longrightarrow Q_X^A.                            \tag{D.78.21}
\]
The forward implication follows from equality of the subgroups.
For the reverse implication, take a nonidentity $k_0\in K$.
The compact subgroup $AT$ is closed, and $k_0\notin AT$:
its $u$-label is zero, so membership would put it in $A$, where $\chi$
is injective. Choose an open normal $N\triangleleft E$ contained in
$F$ such that $k_0\notin ATN$. For the finite set $X=E/(ATN)$,
the identity coset is fixed by both $A$ and $T$, whereas $k_0$ moves it.
The corresponding basis vector in $Q_X$ lies in $Q_X^A$ and not in
$Q_X^B$. Thus (D.78.21) fails for that actual finite coefficient.

The cusp restriction (D.78.10) vanishes on the entire other factor.
The perfect duality proves (D.78.1)--(D.78.3), but has supplied no
surjectivity for (D.78.21). This is the retained coefficient map that
still detects the full pro-odd kernel. The coefficient comparison retains the actual
$K$-orbit partition and supplies no replacement of $Q_X^K$ by $Q_X$. The original cusp $c$, unique transporter $g$, and
D.75 conjugacy equation are unchanged throughout.

\end{proof}
